% ===================================================================
%  GAMBAR No. 10 — Segara. --- Pelabuhan.
%
%  Traduction, faite en 2026, en javanais d'aujourd'hui, sur l'ido de
%  texto/io. Regles de la colonne : voir 10-gambar-01.tex. Une seule
%  numerotation.
%
%  LA MEME PARTITION QU'A LA COLONNE OURDOUE, ET ELLE TOMBE ICI AU
%  MEME ENDROIT. Le releve ourdou avait trouve, au tableau 10, que la
%  mer se partage entre la VOILE, qui garde ses mots d'ici, et la
%  VAPEUR, qui prend ceux d'Europe. Le javanais fait exactement le
%  meme partage, et il est encore plus net parce que Java est une ile.
%
%   * La voile est javanaise et malaise sans un emprunt : prau le
%     bateau (46), layar la voile, dhayung la rame (47), tiang layar
%     le mat (51), jangkar l'ancre (77), geladhak le pont (48), jala
%     le filet (56), tampar le cordage (52), ombak la vague (72),
%     pasang lan surut le flux et le reflux, teluk la baie (95),
%     tanjung le cap, pulo l'ile (99), karang le recif (100), selat le
%     detroit (98), pesisir la plage (73), muara l'embouchure.
%   * Le navire d'acier est europeen d'un bout a l'autre : kapal uap
%     (3), kapal selam (10), torpedo (8), kapal penjelajah (6),
%     fregat (11), dok (26), baling-baling (28), kemudhi (30), derek
%     (32), ketel (25), semafor (70), trem listrik (74).
%
%  ET LE MOT QUI TIENT LES DEUX EST « kapal », qui est tamoul (kappal)
%  arrive par le malais : il nomme le gros navire depuis des siecles,
%  et il a servi tout naturellement pour le cuirasse. Le javanais ne
%  dit pas prau pour un navire de guerre — la difference de taille est
%  dans le lexique avant d'etre dans l'acier.
%
%  ET C'EST ICI QUE LE CINQUIEME MOT SORTI DE JAVA SE TROUVE, APRES
%  قلی, بازار, wajan ET lahar. « jong », le grand voilier javanais, a
%  ete pris par les Portugais sous la forme « junco », et l'anglais en
%  a fait « junk », le nom qu'il donne aujourd'hui au voilier chinois.
%  Un mot javanais est parti d'ici pour revenir designer un bateau qui
%  n'est pas d'ici. C'est le seul des cinq qui ait change de chose en
%  route.
%
%  L'AMIRAL (84) SE DIT « laksamana », ET C'EST LE PLUS BEAU TITRE DU
%  LIVRET. Le mot vient de Lakṣmaṇa, le frere de Rama dans le
%  Ramayana, et il est devenu le titre du chef de flotte des
%  sultanats malais, puis le grade d'amiral de la marine indonesienne
%  et le mot javanais d'aujourd'hui. Le lieutenant (78) et le
%  capitaine (79) sont neerlandais ; seul le plus haut grade est reste
%  au poeme. Ce n'est pas un hasard : c'est le seul rang qui existait
%  avant que le Hollandais arrive.
%
%  ET LE PILOTE (36) SE DIT « pandhu laut », de pandhu, qui est le
%  guide — le meme mot que les Pandhawa du wayang, ceux qu'on guide.
%  Le port est le seul tableau du livret ou le lexique du poeme et
%  celui de la machine se partagent la page ligne a ligne.
% ===================================================================

%%K t10-apar-1 apar
\VUsaut{4.0mm}
\VUtitre{15.0pt}{60}{GAMBAR No. 10}
\VUsaut{3.0mm}
\VUpk{11.4pt}{40}{Segara. --- Pelabuhan.}
\VUsaut{3.0mm}

%%K t10-01-1 p
1. --- Ing mangsa panas, nalika sawetara wong nggoleki tentrem lan hawa
adhem ing gunung, liyane luwih seneng menyang pesisir. Mangkono kang
kita lakoni taun kepungkur, amarga kita seneng banget marang
\VUgras{Samodra} \textsuperscript{(1)}.

%%K t10-02-1 p
2. --- Tumrap aku, aku rumangsa seneng banget nyawang jembare banyu
kang tanpa wates kang megar tekan \VUgras{wates langit}
\textsuperscript{(2)}, lan weruh mangkate \VUgras{kapal uap}
\textsuperscript{(3)} kang gedhe temen kanggo \VUgras{lelayaran} kang
dawa, panggonan para lelungan kang menyang Amerika padha munggah.

%%K t10-03-1 p
3. --- Mula bapakku, kang ngerti senengku, tansah milih, kanggo
nglakoni preinan, pesisir kang tentrem banget, kang manggon cedhak
salah sijine \VUgras{pelabuhan perang} kita kang gedhe.

%%K t10-03-2 p
Nalika kita nginep pungkasan, \VUgras{eskadron kapal} teka labuh ing
mburine \VUgras{tanggul segara} \textsuperscript{(4)} kang gedhe temen,
kang nutup lan ngayomi \VUgras{pelabuhan perang} \textsuperscript{(5)},
lan aku bisa nyawang miturut karepku warna-warna \VUgras{kapal
perang,} \VUgras{kapal selam} \textsuperscript{(10)} cilik kang nyilem
lan ilang ing sangisore banyu, lan uga \VUgras{kapal perang lapis
waja} \textsuperscript{(9)} kang gedhe temen, kang nganti saiki meh mung
wedi marang serangane \VUgras{kapal torpedo} \textsuperscript{(8)}.

%%K t10-tit-2 sub
\VUsaut{3.0mm}
\VUpk{10.2pt}{40}{Kapal-kapal cilik.}
\VUsaut{3.0mm}

%%K t10-04-1 p
4. --- Kapal torpedo iku wis bisa \VUgras{kanthi trampil} nemu titik
ringkihe lapis waja logam kang kandel, kanggo nancepake
\VUgras{torpedo} kang mbebayani, kang bledhosane njalari sirnane kapal
— nanging kapal torpedo isih kurang medeni tinimbang kapal selam,
amarga kapal panulak torpedo gampang mburu dheweke.

%%K t10-05-1 p
5. --- Aku uga bisa weruh \VUgras{kapal penjelajah} \textsuperscript{(6)}
kang cepet lan lapis waja, kang meh padha ora bisa dilarani karo kapal
perang lapis waja kang sanyatane, sawetara \VUgras{kapal angkut}
\textsuperscript{(12)}, telu utawa papat \VUgras{prau meriem}
\textsuperscript{(7)}, lan pungkasane siji \VUgras{fregat}
\textsuperscript{(11)}. \VUgras{Tontonan pepanthan kapal iku, nalika
manuver arep ngunggahake jangkar, narik banget.}

%%K t10-06-1 p
6. --- Sepira bedane wangun antarane tumpukan waja kang gedhe iku karo
\VUgras{kapal tiang telu} kang apik utawa \VUgras{kapal layar}
\textsuperscript{(13)} kang luwes kang nganti saiki isih dianggo ing
pepanthan kapal dagang kang dakdeleng mlebu pelabuhan, kanthi layar
megar, nalika aku mlaku-mlaku ing \VUgras{dermaga}
\textsuperscript{(14)}. Pancen kapal layar iku kelangan akeh kaunggulane
\VUgras{nalika} \VUgras{angin}e nglawan. Kapal-kapal iku kudu ngudhunake
kabeh layare kanggo mlebu maneh menyang kolam, lan \VUgras{kapal
panarik} \textsuperscript{(16)} perlu kanggo njemput lan nggawa
kapal-kapal iku, kaya \VUgras{prau momotan} \textsuperscript{(17)}
lumrah, menyang dermaga.

%%K t10-tit-3 sub
\VUsaut{3.0mm}
\VUpk{10.2pt}{40}{Pelabuhan dagang.}
\VUsaut{3.0mm}

%%K t10-07-1 p
7. --- Kang narik ing pelabuhan kang dakcritakake iku, yaiku yen
pelabuhan kono ora mung isi pelabuhan perang, kang ditata kanthi
pagawean raseksa, nanging uga \VUgras{pelabuhan dagang} kang
nggumunake, kang manggon ing \VUgras{muara} kali gedhe.

%%K t10-08-1 p
8. --- Ilat lemah kang misahake kolam loro iku dibentengi lan dijaga
nganggo urut-urutan \VUgras{baterei meriem} lan \VUgras{bastion} kang
maju menyang segara. Wong butuh idin mligi kanggo mlaku-mlaku ing
\VUgras{tembok beteng} \textsuperscript{(15)}. Aku gampang olehe, liwat
pakdheku, kang dadi insinyur ing \VUgras{papan gawe kapal}
\textsuperscript{(18)}, lan aku lunga nyawang kapal-kapal kang lagi
digawe ing sangisore hangar kang jembar.

%%K t10-09-1 p
9. --- Aku malah nyekseni diluncurake kapal perang lapis waja, lan aku
kelingan wektu kang gawe trenyuh kang dirasakake wong akeh nalika
tumpukan raseksa iku, kang dieculake dhewe, wiwit mlorot ing rangka
kaya ayunan lan banjur ngambang ing banyune kali.

%%K t10-10-1 p
10. --- Kanggo menyang papan gawe kapal, wong ngliwati \VUgras{lapangan
latihan} \textsuperscript{(19)} panggonan para prajurit garnisun saben
dina latihan, banjur liwat \VUgras{kreteg cilik} \textsuperscript{(20)}
wesi kang nyambungake lapangan parade karo \VUgras{gudhang gegaman}
\textsuperscript{(21)}.

%%K t10-tit-4 sub
\VUsaut{3.0mm}
\VUpk{10.2pt}{40}{Gudhang gegaman lan dok.}
\VUsaut{3.0mm}

%%K t10-11-1 p
11. --- Karo pakdheku, aku nyowani bangunan gedhe iku, kang kebak
\VUgras{meriem} \textsuperscript{(22)}, \VUgras{mimis gedhe}
\textsuperscript{(23)} lan \VUgras{obus.} Aku uga weruh \VUgras{pabrik
logam} \textsuperscript{(24)} panggonan digawe \VUgras{ketel}
\textsuperscript{(25)} kapal uap.

%%K t10-12-1 p
12. --- Cedhak banget, ana \VUgras{dok} \textsuperscript{(26)} — dok
ndandani — lan \VUgras{dok ngambang} \textsuperscript{(27)} kang
nggumunake temen, panggonan aku bisa weruh saka cedhak, ing kapal kang
lagi didandani, \VUgras{baling-baling} \textsuperscript{(28)},
\VUgras{lunas} \textsuperscript{(29)} lan \VUgras{kemudhi}
\textsuperscript{(30)} prau.

%%K t10-13-1 p
13. --- Pelabuhan dagang iku padha pantese disinaoni karo pelabuhan
perang, lan aku nglampahi pirang-pirang jam nyawang ing dermaga
panggonan ditambatake \VUgras{kapal rodha} \textsuperscript{(31)} kang
mudhik ing kali, lan nyawang mlakune \VUgras{derek tiang}
\textsuperscript{(32)} kang ngangkat momotan gedhe temen kaya wulu.

%%K t10-tit-5 sub
\VUsaut{4.0mm}
\VUpk{10.2pt}{40}{Pandhu laut.}
\VUsaut{3.0mm}

%%K t10-14-1 p
14. --- Aku gumun marang \VUgras{kapal trans-Atlantik}
\textsuperscript{(34)} kang metu saka pelabuhan, karo \VUgras{gendera}ne
\textsuperscript{(35)} megar ing angin, dituntun \VUgras{pandhu laut}
\textsuperscript{(36)} kang trampil, kang nggumunake anggone bisa
ngetutake \VUgras{alur} kang ditengeri \VUgras{pelampung tandha}
\textsuperscript{(40)} lan \VUgras{cagak tandha,} karo nyingkiri
\VUgras{prau tongkang} \textsuperscript{(41)} kang rekasa olehe dikemudhi
mung nganggo layar pesagine kang siji. Aku bisa mbedakake ing
\VUgras{cucuk kapal} \textsuperscript{(37)} \VUgras{kamar}
\textsuperscript{(39)} kelas loro, lan ing \VUgras{buritan}
\textsuperscript{(38)} ruwang tamu kanggo penumpang kelas siji.

%%K t10-15-1 p
15. --- Kadhang aku uga nyekseni diisenane \VUgras{prau momotan}
\textsuperscript{(42)} nganggo barang kang lagi wae ditumpuk saka
\VUgras{gudhang} \textsuperscript{(43)} lan \VUgras{setasiun laut}
\textsuperscript{(44)}, kang digawa sepur akeh banget ing \VUgras{sepur
dermaga} \textsuperscript{(45)}.

%%K t10-tit-6 sub
\VUsaut{3.0mm}
\VUpk{10.2pt}{40}{Ndhayung kanggo seneng-seneng.}
\VUsaut{3.0mm}

%%K t10-16-1 p
16. --- Aku kerep ndhayung ing \VUgras{dok banyu} \textsuperscript{(53)}
panggonan \VUgras{prau} \textsuperscript{(46)} padha ngeyup, lan bungah
temen rasaku olehe ngobahake \VUgras{dhayung} \textsuperscript{(47)}. Aku
munggah ing \VUgras{geladhak} \textsuperscript{(48)} \VUgras{prau layar}
\textsuperscript{(49)}, aku menek nganggo \VUgras{tampar}
\textsuperscript{(52)} ing \VUgras{tiang layar} \textsuperscript{(51)}
tekan \VUgras{palang layar} \textsuperscript{(50)}, banjur aku mbaleni
lelungan \VUgras{nganggo prau balap} \textsuperscript{(54)} ing antarane
\VUgras{prau mancing} \textsuperscript{(55)} kang abot, panggonan
\VUgras{jala} \textsuperscript{(56)} padha dipepe.

%%K t10-17-1 p
17. --- Ing \VUgras{dermaga} \textsuperscript{(58)} liwat ing ngarepku
\VUgras{barang dagangan} \textsuperscript{(59)} kang maneka warna, kang
diwadhahi \VUgras{peti} \textsuperscript{(60)} utawa \VUgras{buntelan
gedhe} \textsuperscript{(61)}. Kabeh iku dadi \VUgras{momotan}
\textsuperscript{(63)} prau gedhe, lan \VUgras{derek}
\textsuperscript{(62)} kang kuwat njupuk lan nyelehake ing
\VUgras{trek} \textsuperscript{(64)}, dene \VUgras{tukang angkutan}
ngumumake barange lan mbayar pajeg ing \VUgras{pabean}
\textsuperscript{(65)}.

%%K t10-18-1 p
18. --- Pungkasane, nalika aku tekan \VUgras{kreteg puter}
\textsuperscript{(66)} utawa \VUgras{pintu banyu} \textsuperscript{(69)},
liwat ing ngarepe \VUgras{mesin keruk} \textsuperscript{(68)} kang
ngresiki \VUgras{kanal} \textsuperscript{(67)}, aku mudhun ing dermaga
cedhak \VUgras{semafor} \textsuperscript{(70)}.

%%K t10-19-1 p
19. --- Ing sisih sijine dermaga iku pancen \VUgras{sekoci}
\textsuperscript{(71)} padha ngrapet, kang nggawa para opsir eskadron
menyang dharatan. Nggumunake temen nonton para wong laut ngangkat
dhayunge kanthi ajeg, dene prau kang digawa \VUgras{ombak}
\textsuperscript{(72)} kanthi gampang ngrapet ing undhak-undhakane papan
mudhun. Ing panggonan iku, wong bisa luwih becik nyawang
\VUgras{pasang} lan obahe \VUgras{pasang} lan \VUgras{surut} ing
\VUgras{pesisir} \textsuperscript{(73)}.

%%K t10-tit-7 sub
\VUsaut{3.0mm}
\VUpk{10.2pt}{40}{Papan kumpulan.}
\VUsaut{3.0mm}

%%K t10-20-1 p
20. --- Kanggo bali menyang omah, aku numpak \VUgras{trem listrik}
\textsuperscript{(74)} kang \VUgras{halte}ne \textsuperscript{(75)} ana
ing cedhake \VUgras{cagak} kang dipasang ing ngarepe papan kumpulane
para opsir.

%%K t10-20-2 p
Saka \VUgras{balkon} \textsuperscript{(76)} papan kumpulan iku, kang
direngga \VUgras{jangkar} \textsuperscript{(77)}, meriem lan mimis
gedhe, aku kerep weruh kumpulan endah para opsir laut :
\VUgras{letnan} \textsuperscript{(78)} karo pedhange, \VUgras{kapten
artileri} \textsuperscript{(79)} lan \VUgras{kapten infanteri}
\textsuperscript{(80)} karo \VUgras{pedhang}e. Ing mburine ana
\VUgras{kelasi} \textsuperscript{(81)} kang nggawakake \VUgras{teropong}
yen para opsir kepengin mbedakake apa kang kedadeyan ing kadohan.

%%K t10-21-1 p
21. --- Aku uga weruh \VUgras{letnan enom} \textsuperscript{(82)} nampa
prentahe \VUgras{komandan} \textsuperscript{(83)} utawa
\VUgras{laksamana} \textsuperscript{(84)}, dene \VUgras{kelasi kepala}
\textsuperscript{(85)} ngenteni kanggo nggawa prentah iku menyang kapal.

%%K t10-22-1 p
22. --- Saka balkon iku, pemandangane endah banget : wong nguwasani
kabeh pesisir karo \VUgras{tarub}e \textsuperscript{(86)} lan
\VUgras{kamar salin}e \textsuperscript{(87)}, lan wong weruh, ing jame
adus, \VUgras{wong-wong kang adus} \textsuperscript{(88)} kang
pethakilan bungah ing ngarepe \VUgras{kasino} \textsuperscript{(89)}.

%%K t10-23-1 p
23. --- Ing pucuke pesisir, dharatan iku mbentuk \VUgras{semenanjung}
\textsuperscript{(91)} cilik kang \VUgras{akeh watune,} panggonan
\VUgras{ombak gedhe} \textsuperscript{(92)} teka pecah, lan ing kono
dibangun \VUgras{mercusuar} \textsuperscript{(93)} kang ing wayah bengi
nuduhake dalan marang para kang lelayaran. Semenanjung iku sesambungan
karo dharatan mung liwat \VUgras{tanah genthing} kang ciyut banget
\textsuperscript{(90)}.

%%K t10-tit-8 sub
\VUsaut{3.0mm}
\VUpk{10.2pt}{40}{Pemandangan umum pesisir.}
\VUsaut{3.0mm}

%%K t10-24-1 p
24. --- \VUgras{Tebing curam} \textsuperscript{(94)} megar dawa, kabelah
dening segara, kang kanthi sakarepe nggambar ing kono urut-urutan
\VUgras{teluk} \textsuperscript{(95)} lan \VUgras{tanjung,} kang gawe
tebing iku endah banget.

%%K t10-24-2 p
Ing \VUgras{dhataran dhuwur} \textsuperscript{(96)}, kang nguwasani
\VUgras{selat} \textsuperscript{(98)}, ana \VUgras{beteng}
\textsuperscript{(97)} kang wigati banget lan sawetara « vila ».

%%K t10-25-1 p
25. --- Selat iku misahake saka bawana \VUgras{pulo}
\textsuperscript{(99)} kang mbebayani banget, kang kinubeng
\VUgras{karang} \textsuperscript{(100)} lan \VUgras{ombak pecah,}
panggonan prau lan malah kapal gedhe kadhang kandhas. Sanajan
pitulungan cepet teka, lan \VUgras{prau pitulungan} enggal tekan,
\VUgras{karame kapal} iku medeni banget, lan, sajrone \VUgras{prahara}
ekuinoks, pirang-pirang kapal sirna bareng wong lan barange.

%%K t10-25-2 p
Sanajan cedhak karo bebaya iku, \VUgras{pelabuhan} iku rame banget
diparani para wong laut.

\VUsaut{6.0mm}
\VUfilet{22mm}
